\section{Study Design and Tests}
\label{sec:methods}

A \emph{behavioral profile} estimates how often a model performs an action.
A default profile describes its usual rate; a conditional profile also
accounts for its responses to student cues. We examine rated teaching
behaviors, expressed confidence, safety boundaries, affiliation, and specific
actions. These analyses use different model groups and measures
(Appendix~\ref{app:behavior-evidence}); their samples overlap and cannot be
counted as independent observations.

\paragraph{What stability means.}
We ask whether a model repeats an action on the same input, whether models
show similar relative tendencies across tasks, and whether past behavior helps
predict new responses. These are separate questions. Showing that action rates
remain unchanged requires an equivalence test. Different forms of the intraclass
correlation coefficient (ICC) assess agreement among annotators and consistency
of model differences across tasks.

\subsection{Historical Data and Controlled Experiments}
Historical benchmark responses support exploratory comparisons of teaching
behavior, confidence, safety, and responses to instructions. Each analysis
uses a relevant subset; Appendices~\ref{app:behavior-evidence}
and~\ref{app:archive-scope} describe model coverage and data limitations.

Controlled experiments test five models: MiniMax-M3, MiniMax-M2.7,
GLM-5.3, DeepSeek-V4-Pro, and Doubao-Seed-2.0-Lite. Here, a model refers to the
version and settings tested. We use 32 training and 32 test questions from
separate groups of related problems, balanced across mathematics, science,
language, and logic/data reasoning. Research agents authored the English
problems and student messages; two model reviewers checked them before the
materials were finalized. The tests use new problems within these domains.

For each problem, we combine two student-work states (erroneous or correct but
unfinished) with two statements of feeling (calm or frustrated). Tutors receive
the same reference solution, unseen by the student. We request two responses
per input and group comparisons by source problem. Appendix~\ref{app:detailed-methods}
records materials, grouping, the separate measurement pilot, and model settings.

\subsection{Observed Actions and Repeated Responses (RQ1)}
We focus on three actions that can be identified in the text and for which
annotators agreed on positive examples: \emph{answer revelation}, an explicit
final answer regardless of correctness; \emph{reasoning elicitation}, an
invitation for student reasoning not immediately answered by the tutor; and
\emph{affect acknowledgement}, explicit recognition of feeling or difficulty
beyond generic praise or politeness. A reply can contain several actions;
these are not assumed to be independent personality factors.
Five secondary actions are retained regardless of results.

Three models---MiniMax-M3, DeepSeek-V4-Pro, and GLM-5.3---label the actions in
each response using its context, reference, and numbered answer lines.
These annotators are not given the generating model's name or teaching
instruction, although the answer itself may offer clues. Each positive label
must cite supporting lines. A majority vote determines each action label.
We retry or repair only incomplete or incorrectly formatted annotations;
disagreement or label values do not trigger recovery
(Appendix~\ref{app:detailed-methods}).
RQ1 reports action rates and agreement between repeated responses, separating
agreement on presence from agreement on absence. Statistical comparisons
account for responses sharing a source problem.

\subsection{Predicting Responses to New Problems (RQ2)}
We use logistic regression to predict each action's probability for the same
five models on new questions. We fit predictors to labeled training responses
using model identity, subject, and scripted student state, and fix all
predictions before collecting test responses.

\emph{Default gain} measures how much knowing a model's usual behavior improves
prediction beyond subject and student state. \emph{State gain} measures the
added value of knowing how that model responds to student cues, beyond its
usual behavior within each subject. Both comparisons already account for
student state; the second adds model-specific responses to it.
Both gains are reductions in Brier loss, a measure of probability prediction
error. Positive gains mean more accurate predictions.
Appendix~\ref{app:prediction-method} gives fitting and statistical testing details.

\subsection{Changing How Students Express Their State (RQ2)}
\label{sec:cue-design}
We change how students express their state and collect new responses from the
same five models on the same 32 test problems. We then test whether the earlier
predictions still work, without refitting the predictors. This tests new wording
on the same test problems. We fixed this follow-up design after seeing the
first test results, but before collecting the new responses.

We compare original wording rerun during the same period, new explicit
expressions, implicit cues, and a classmate control: 5,120 additional responses
and 10,440 across both stages, excluding the pilot. Explicit and implicit
conditions each use four calm/frustrated cue pairs, balanced across subjects.
Implicit cues also change actions, willingness, or time information, so they
do not isolate emotion. In the classmate control, the current student remains
calm while quoting a calm or frustrated classmate working on another problem.

Six primary tests examine default gains for giving answers and inviting
reasoning, and state gains for acknowledgement, under explicit and implicit
wording, with Holm correction. Original-wording and classmate comparisons
are secondary.
Appendix~\ref{app:cue-transfer} gives exact cues, uncertainty estimates,
annotation recovery, and checks excluding individual annotators and models.

\subsection{Instruction Effects and Robustness Checks (RQ3)}
On 16 problems selected by a fixed hash rule, each model receives every student
state under the standard tutoring prompt, two neutral paraphrases, an
instruction to invite reasoning without the answer (ASK), and one to explain
and give the answer (EXPLAIN). Paired comparisons measure changes in action
rates and differences in these changes across models and student states.
With an equivalence margin of $\pm0.10$, the sample is too small to establish
equivalent effects of neutral paraphrases beyond these problems. We also
measure how often answers and reasoning invitations appear together, without
inferring their order or teaching quality.

Checks specified before test responses were collected use eight sets of
annotations, exclude each generating model in turn, and resample training
problems 200 times. To check sensitivity to answer errors, we exclude matched
response groups from all five models without refitting predictions. Reviewers'
low agreement on errors limits this check. Appendix~\ref{app:detailed-methods}
gives the full procedures, diagnostic controls, and protocol changes.

A separate analysis of historical MathTutorBench responses
\citep{macina2025mathtutorbench} compares general scaffolding and explicit
probing, both limited to two sentences, on 256 matched dialogue groups from
seven models. We test whether each model's behavior under one instruction
predicts its responses under the same or the other instruction, compared with
knowing only the current instruction. Two annotations remain missing because
the provider blocked the requests; the other annotators agree, determining
each majority label. We also analyze the 254 fully annotated groups
(Appendix~\ref{app:archive-validation}).
